\documentclass[10pt,a4paper]{article}
\usepackage[utf8]{inputenc}
\usepackage{amsmath}
\usepackage{amsfonts}
\usepackage{amssymb}
\usepackage{graphicx}
\usepackage{float}
\usepackage{booktabs}
\usepackage{xcolor, soul}
\sethlcolor{green}
\usepackage{enumerate}
\title{GSF-3100 \\
07 Marché des obligations corporatives \\
Série d'exercices 7 (Solutions)}

\begin{document}
\maketitle

\begin{itemize}
\item Dans ce document vous avez des questions à choix multiples allant de a) à d).  
\item Il y a seulement un choix possible par question.
\item Ce questionnaire contient des exercices en lien avec le marché des obligations corporatives et plus pécisément sur la prime de risque et la prime d'insolvabilité.
\item Contient 3 questions. 
\end{itemize}
 
\newpage

\section*{Mise en contexte}

Un investisseur considère une obligation corporative venant à échéance dans 3 périodes avec un taux de coupon de 2 \% et un taux de rendement à l’échéance de 1,5 \% (effectif). Il prévoit les scénarios d’insolvabilité suivants (avec leur probabilité entre parenthèses): 

\begin{enumerate}
\item Aucun problème (50 \%): Les coupons et la valeur nominale sont tous remboursés à temps;

\item Difficulté financière (25 \%): Le versement des coupons est suspendu et seule la valeur nominale est remboursée à l'échéance;

\item Faillite (25 \%): Seul le premier coupon est versé à l'investisseur. La valeur nominale n'est jamais remboursée.
\end{enumerate}

\section*{Exercice 1}
Trouvez le taux de rendement espéré de l'obligation corporative.

\begin{enumerate}[a)]
\item 22.5\%
\item 23.9\%
\item -22.5\%
\item \hl{23.9\%}
\end{enumerate}


\section*{Exercice 2}
Trouvez la prime d’insolvabilité de l'obligation corporative.
\begin{enumerate}[a)]
\item 23.4\%
\item 24.4\%
\item \hl{25.4\%}
\item 26.4\%
\end{enumerate}

\section*{Exercice 3}
Trouvez la prime de risque de crédit de l'obligation corporative,  en supposant que l’obligation gouvernementale de référence offre un taux de rendement à l’échéance (effectif) de 8 \%
\begin{enumerate}[a)]
\item 6.5\%
\item \hl{-6.5\%}
\item 8.5\%
\item -8.5\%
\end{enumerate}

\end{document}